\bigskip
\noindent Autres pistes de l'onglet Findings pour ce thème, dans des cotes qui n'ont pas de section ici :\par

\begin{encadre}
\begin{theoreme*}[Grothendieck, pages 144 et 145, quotients des feuilletages ; voir aussi Kock--Moerdijk~\cite{kockmoerdijk1992}\piste{119-foliation-quotient-topos-equivalence}]\nonverifie
La catégorie des couples $(X, R)$, où $X$ est un espace topologique et $R$ une relation d'équivalence locale sur $X$, localement ouverte et à fibres localement connexes, est équivalente à celle des triplets $(X, \mathcal Y, f)$, où $\mathcal Y$ est une \q{multiplicité topologique} et $f : \mathrm{Top}(X) \to \mathcal Y$ un morphisme localement ouvert à fibres localement connexes.
\end{theoreme*}
\end{encadre}
{\small La page l'énonce sans démonstration. L'édition rapproche les multiplicités topologiques des étendues, topos localement équivalents à des topos de faisceaux sur un espace, et les relations d'équivalence locales de celle des feuilles d'un feuilletage. Kock et Moerdijk (1992) \cite{kockmoerdijk1992} montrent que toute étendue est équivalente au topos $\mathrm{sh}(L, r)$ des $r$-faisceaux sur un locale $L$ muni d'une relation d'équivalence locale $r$ localement simplement connexe, et que $\mathrm{sh}(L, r)$ est une étendue pour un tel $r$ ; leur introduction renvoie à SGA~4, exposé~IV, pour les relations d'équivalence locales. Kock et Moerdijk (1996) \cite{kockmoerdijk1996} identifient $\mathrm{sh}(M, r)$ aux faisceaux sur un groupoïde étale sous des hypothèses de connexité. L'équivalence de catégories de la page, entre couples $(X, R)$ et morphismes $\mathrm{Top}(X) \to \mathcal Y$, n'a pas été vue dans les pages lues de ces deux articles. Voir aussi les morphismes de topos des sections~\ref{sec:19} et~\ref{sec:161-1}. Mathlib n'a pas les relations d'équivalence locales ni les multiplicités topologiques. On se reportera aux notes de Grothendieck (\lot{119}{8}{144}).\par}
